\section{常返、暂留与返回时间}
\label{sec:06-Random-Processes/03-Markov-Chains/03}

做用户留存的时候你迟早会碰到这个问题：一个本周活跃的用户，沉寂了三周，还算不算流失？

把它写成链。状态里有一个叫``活跃''，用户在活跃与各种不活跃状态之间游走。问题变成：从活跃出发，最终回到活跃的概率是不是一？如果是一，那么无论沉寂多久，他都还会回来，而且以概率一回来无穷多次——所谓``流失''只是等待时间长而已。如果那个概率是 \(0.9\)，情况完全不同：每次回来之后都要再赌一次，早晚有一次赌输，此后永不再见。

上一节的可达性只保证存在一条通向``活跃''的正概率路径。路径存在，不等于路径会被走上。本节要分的就是这条界线。

\subsection{概率恰好等于一，才叫常返}

从状态 \(i\) 出发，定义首次返回时刻

\[
T_i^+=\inf\set{n\ge1:X_n=i}.
\]

下标从 \(n\ge1\) 起算，不把``出发时就在 \(i\)''那一刻算进去；这个细节后面会反复用到。概率符号里的下标表示初始条件，\(P_i(\cdot)=P(\cdot\given X_0=i)\)，期望 \(\E_i\) 同理。

若 \(P_i(T_i^+<\infty)=1\)，称 \(i\) 为常返状态；若这个概率严格小于一，称暂留状态——有正概率一去不返。这里没有中间地带，判据卡在概率等不等于一这一个点上，而不是等不等于某个较大的数。

\subsection{回来一次，就等于回来无穷多次}

设 \(i\) 常返。链第一次回到 \(i\) 的那个时刻是一个停时，在那一刻链重新站在 \(i\) 上，且强 Markov 性保证之后的行为与最初从 \(i\) 出发时同分布。于是``再返回一次''的概率仍然是一。归纳下去，对任意 \(k\)，返回至少 \(k\) 次的概率都是一；取极限，常返状态以概率一被访问无穷多次。

暂留状态走的是同一条论证，结论相反。记 \(\rho=P_i(T_i^+<\infty)<1\)。每一次访问结束后，链以概率 \(\rho\) 还会回来，以概率 \(1-\rho\) 永远离开，而且由 Markov 性，每次的赌局互相独立、赔率相同。访问次数因此服从几何分布，期望为

\[
1+\rho+\rho^2+\cdots=\frac1{1-\rho}<\infty.
\]

\begin{center}
\begin{tikzpicture}[x=1cm,y=1cm,font=\small,>=stealth]
  \foreach \k in {0,1,2,3}{
    \draw[black!55,fill=blue!8] (2.1*\k,0) circle (0.42);
  }
  \node at (0,0) {\(i\)};
  \node at (2.1,0) {\(i\)};
  \node at (4.2,0) {\(i\)};
  \node at (6.3,0) {\(i\)};
  \foreach \k in {0,1,2}{
    \draw[->,thick] (2.1*\k+0.46,0) -- (2.1*\k+1.64,0);
    \node[font=\footnotesize,above] at (2.1*\k+1.05,0.08) {\(\rho\)};
  }
  \draw[->,thick,black!50] (6.76,0) -- (7.9,0);
  \node[font=\footnotesize] at (8.3,0) {\(\cdots\)};
  \foreach \k in {0,1,2,3}{
    \draw[->,red!65] (2.1*\k,-0.46) -- (2.1*\k,-1.3);
    \node[font=\footnotesize,right,red!65] at (2.1*\k+0.08,-0.9) {\(1-\rho\)};
  }
  \draw[red!45,fill=red!6] (-1.0,-2.1) rectangle (7.3,-1.35);
  \node[font=\footnotesize,black!70] at (3.15,-1.72) {永远离开，不再回来};
  \node[font=\footnotesize,black!60,align=center] at (-1.2,0.95) {第 \(1,2,3,\ldots\) 次\\访问 \(i\)};
\end{tikzpicture}

\emph{每一次身处 \(i\)，链都要重新掷一次同样的硬币。\(\rho=1\) 时红色的那条出口不存在，访问次数无穷；\(\rho<1\) 时出口每次都开着，迟早会被走掉，访问次数是几何分布，期望 \(1/(1-\rho)\)。}
\end{center}

暂留状态可以被访问很多次——\(\rho=0.99\) 时平均一百次——但总归是有限次。常返与暂留的差别不是``多''与``少''，是有限与无限。

\subsection{一个级数就能判定}

直接验算 \(P_i(T_i^+<\infty)=1\) 要对无穷多条路径求和，很难下手。上面那张图给出了一条捷径。\((P^n)_{ii}\) 是第 \(n\) 步恰好落在 \(i\) 的概率，对 \(n\) 全部加起来，正是从 \(i\) 出发访问 \(i\) 的期望总次数：

\[
\sum_{n=0}^{\infty}(P^n)_{ii}
=\E_i[\text{访问 }i\text{ 的总次数}].
\]

于是判据自动浮现：级数发散则期望次数无穷，\(i\) 常返；级数收敛则期望次数有限，\(i\) 暂留。你不需要知道 \(\rho\) 的数值，只要判断一个正项级数的敛散。

\subsection{常返性是通信类的标签}

若 \(i\leftrightarrow j\)，两者要么都常返，要么都暂留。骨架只有一行：取 \(m,n\) 使 \((P^m)_{ij}>0\) 且 \((P^n)_{ji}>0\)，则对任意 \(k\)，

\[
(P^{m+k+n})_{jj}\ge(P^n)_{ji}\,(P^k)_{ii}\,(P^m)_{ij},
\]

因为右边只是绕经 \(i\) 的那一族路径，左边是所有路径。于是 \(\sum_k(P^k)_{jj}\) 被 \(\sum_k(P^k)_{ii}\) 乘一个正常数控制住；把 \(i\) 与 \(j\) 对调再来一遍，两个级数同敛散。

所以常返性不是单个状态的私事，而是整个通信类的标签，这与上一节的周期一样。有限链上还能说得更死：每个闭通信类都常返，因为链被关在一个有限的集合里，不可能对每个状态都只访问有限次；不在任何闭类里的状态则必然暂留，链早晚会掉进某个闭类再也出不来——上一节那张图里的 \(C\) 就是例子。

``有限''这个词在上一段里承担了全部重量。无限状态的不可约链可以全体暂留：空间足够大，链有正概率一路漂走。

\subsection{回来了，但平均要等多久}

常返还能再分两层。正常返要求 \(\E_i[T_i^+]<\infty\)，不但一定回来，平均等待时间还有限；零常返则是 \(\E_i[T_i^+]=\infty\)，以概率一回来，可期望被返回时间分布的重尾拖散了。

整数线上的简单对称随机游走是零常返的标准例子：它以概率一回到原点，平均返回时间却是无穷。这不是悖论——概率一说的是``每条路径终将回来''，期望无穷说的是``回来得慢的那些路径太重''。两句话都对，回答的是不同的问题。

维数在这里制造了一个漂亮的相变。在 \(\Z^d\) 格点上每步等概率走向相邻格点，\(d=1\) 与 \(d=2\) 时原点常返，\(d\ge3\) 时原点暂留——空间大到游走者有正概率永远找不到回家的路。流传的说法是醉汉总能摸回家，醉酒的信鸽却可能永远迷失在天上。这个结果值得记住的地方在于：每一步都无偏，离``最终一定回来''还差着一段，而这段距离由空间的维数决定。\hyperref[sec:06-Random-Processes/04-Applications/01]{``随机游走：从局部步进到整体尺度''}会把击中概率与尺度关系讲得更细。

\subsection{Kac 公式把返回时间翻译成平稳概率}

不可约正常返链有唯一平稳分布 \(\pi\)，并且

\[
\E_i[T_i^+]=\frac1{\pi_i}.
\]

这个等式值得盯一会。它说平稳分布不是解方程解出来的一组抽象数字，而是返回频率的倒数：你在状态 \(i\) 上花的时间占比是 \(0.2\)，就意味着平均每五步回一次 \(i\)。两个看起来分属不同问题的量——``长期占比''和``等多久''——原来是同一件事的两种读法。

平均每五步不等于每五步，返回时间的方差可以很大；但趋势是硬的：占比越高的状态，回访越密。反过来看，若一条无限链的状态全部零常返，平均返回时间无穷，Kac 公式暗示 \(\pi_i\) 只能是零——这时可能还存在满足 \(\mu=\mu P\) 的不变测度，但它的总质量是无穷，归一化不成概率分布。

最后把几个容易混用的词分开。返回时间要求至少走一步再回到出发点；击中时间是从某处首次进入一个目标集合，允许零步（出发就在目标里则击中时间为零）；访问次数数的是被经过的总次数；停留时间则是连续时间链里每次进入某状态后待着的时长。定义里 \(n\ge1\) 与 \(n\ge0\) 的一字之差不是形式主义——若返回时间也允许零步，那么每个状态都零步``返回''，整个分类就塌了。

常返性回答了链会不会回来，周期性回答了它踩什么节拍回来。下一节把这两条线索与平稳方程拧在一起，回答长期分布存不存在、唯不唯一、以及从任意初值出发到底收不收敛。
